% =====================================================================================
% 致谢页 / Acknowledgement page  (mandatory, 1-2 pages, 500-1500 字)
%
% INSTRUCTIONS FOR THE STUDENTS -- read this before editing:
%   * Every bracketed [待填写] item is a fact only you can state truthfully. The competition
%     disqualifies undisclosed paid guidance and incomplete AI disclosure, so fill these in
%     rather than softening them.
%   * The AI-use paragraph below is already accurate as written (it matches
%     research/AI_ASSISTANCE_LOG.md). Keep it, or rewrite it in your own words, but do not
%     delete the tool name, version, stages, purpose, time or frequency, which the rules
%     require explicitly.
%   * Length: this skeleton runs to about 900 characters of Chinese once the brackets are
%     filled; the requirement is 500-1500 字.
% =====================================================================================

\newpage
\section*{致谢与人工智能使用声明\\\large Acknowledgement and Disclosure of AI Use}
\addcontentsline{toc}{section}{致谢与人工智能使用声明 (Acknowledgement / AI disclosure)}

\subsection*{一、研究背景}

本项目研究编码智能体（coding agent）的失败方式，以及运行时系统能否在失败发生\emph{之前}判断智能体
究竟是「找不到代码」（LOST）还是「找到了代码但修改不正确」（WRONG-FIX）——因为这两种情况需要\emph{相反}
的干预措施。研究使用的数据全部来自公开的智能体运行轨迹数据集（SWE-agent / SWE-rebench、
Terminal-Bench 2.0，以及另外四个开源脚手架），以及八个开源 Python 软件包在本机运行的测试套件。
不涉及人类受试者，不使用任何非公开数据。

\subsection*{二、指导老师与参赛学生的关系}

[待填写] 需明确说明以下四点，缺一不可：
\begin{enumerate}[leftmargin=1.6em,itemsep=1pt]
  \item 指导老师的姓名与所在单位；
  \item 学生通过什么渠道认识指导老师（校内课程、科创项目、亲友介绍等）；
  \item \textbf{指导是有偿还是无偿}（如为有偿，须写明费用与来源；如为无偿，亦须明确写出）；
  \item 指导老师具体参与与未参与的部分（例如：确定选题方向、审阅稿件、提供算力；未参与数据采集、
        未参与代码编写等）。
\end{enumerate}
据实填写即可。组委会明确将「未申报的有偿指导」列为取消资格的情形，因此此处不需要顾虑措辞。

\subsection*{三、分工说明}

按组委会要求，逐项说明。括号内为我方建议的填写格式，请按实际修改：

\begin{itemize}[leftmargin=1.3em,itemsep=2pt]
  \item \textbf{选题来源}：\ [待填写] 例如：由 XXX 提出，经与老师讨论后确定为「用运行时路由器区分两
        种失败模式」。
  \item \textbf{数据获取}：\ [待填写] 例如：XXX 负责检索并下载六个公开数据集；XXX 负责搭建
        Python 环境与测试套件。
  \item \textbf{分析与计算}：\ [待填写] 例如：统计分析与模型训练由 AI 助手完成代码编写，XXX 与
        XXX 负责核验每一项结论是否与其产物一致，并推翻了其中十项错误结论（见下）。
  \item \textbf{实验设计与实施}：\ [待填写] 例如：活体实验（145 次真实运行，花费 1.36 美元）由
        XXX 设计三种条件（提示位置／不提示／遮蔽失败位置），由 AI 助手执行运行脚本。
  \item \textbf{撰写论文}：\ [待填写] 例如：论文结构由 XXX 确定，初稿由 AI 助手生成，最终文字由
        XXX 与 XXX 重写并逐句核对。
  \item \textbf{遇到的困难及解决经过}：\ [待填写] 例如：最初使用的定位指标是「自指」的，导致结论
        完全反向；用独立的官方补丁重新计算后方向反转，并在两个留出数据集上复现。
\end{itemize}

\subsection*{四、人工智能使用情况}

\textbf{工具与版本}：\texttt{deepseek-v4.1-flash}（大语言模型），通过 DeepSeek Harness（DSH）
自主编码助手接口调用。每一次模型调用的时间戳、token 数与费用均记录在
\texttt{research/results/rebuild/llm\_ledger.jsonl} 中；活体实验（145 次运行、42 个任务）共花费
\textbf{1.36 美元}。

\textbf{使用阶段与用途}：（1）文献检索与核实（每条关键引文均由人工回原文献核对，其中两条前期结论
因无法查证而被剔除）；（2）编写轨迹解析与数据表构建代码；（3）编写并运行统计分析（任务不交叉的交叉
验证、自助法区间、Wilcoxon 检验、校准曲线、决策曲线）；（4）在活体实验中，模型本身是被测试的\emph{对象}；
（5）编写初稿文字，由参赛学生审阅并重写。

\textbf{时间与频率}：\textbf{2026 年 9 月 10 日至 9 月 13 日}，共四个工作时段，逐次记录在
\texttt{research/AI\_ASSISTANCE\_LOG.md}（共 32 次记录）。

\textbf{聊天记录}：按组委会要求另行提交。

\subsection*{五、关于 AI 出错的自愿声明}

为如实说明研究的可信度，我们列出 AI 助手自身被推翻的十项结论，包括：一个自指的定位指标使结论
方向相反；一项「零误报」声称属于循环论证；一项基线实际上等于运行长度；一次活体实验出现
「48 次全部失败」的惊人结果，实为解释器被删除的实验故障；以及论文中一张移植结果表引用了错误的
数据行——后者由自动审计脚本发现。这说明本项目的规则是：\emph{只有当试图推翻某结论的检验失败时，
该结论才被保留}。

\subsection*{六、声明}

我们声明：本文所述研究成果为参赛学生的原创工作；引用他人成果均已注明并在参考文献中列出；上述
人工智能使用情况完整、真实；我们已知悉并遵守本届竞赛的学术诚信规定。

\vspace{1.2em}
\noindent
\begin{tabular}{ll}
参赛学生签名： & \underline{\hspace{4.5cm}} \quad \underline{\hspace{4.5cm}} \\[10pt]
              & \hfill Ziheng Yu \hspace{2.2cm} Xuhao Chen \\[14pt]
指导老师签名： & \underline{\hspace{4.5cm}} \\[14pt]
日期：        & \underline{\hspace{4.5cm}} \\
\end{tabular}
